\documentclass[a4paper, 12pt]{article}

\usepackage{utils}

\renewcommand*{\today}{02 February 2026}

\begin{document}

\hotbox{Questions ouvertes}{CM 3}{\today}

\section{\today}

\subsection{Ce qui à été fait}

\begin{definition}
    Soit $n, m \in \N$ avec $m \neq 0$.
    Un nombre $n$ est dit $m$-tournant si \text{$n \times m = d10^{l - 1} + \dfrac{n - d}{10}$} où $d$ est le chiffre des unités de $n$ et $l$ sa taille (nombre de chiffres).
\end{definition}

\begin{exemple}
    $$
    128205 \times 4 = 512820
    $$
\end{exemple}

La longueur minimum d'un nombre $m$-tournant est égale à l'ordre de $10$ dans $\Z/(\dfrac{10m - 1}{pgcd(d, 10m - 1)})\Z$.

Voir annexe B la démonstration

\subsection{Les nouvelles questions}

Est-ce possible de généraliser à une base $b$ quelconque ? Si oui, comment trouver la longueur d'un nombre tournant à partir de son multipliant $m$, de son chiffre des unités $d$ et de sa base $b$ ?

Est-ce que pour toute longueur il existe un nombre tournant de cette longueur ?

\end{document}